\documentclass[Master]{subfiles}
\begin{document}
\begin{abstract}
This Master Thesis demonstrates that the non-abelian exchange of Majorana modes in the 2D chiral Topological p-Wave Superconductor vortices is possible with a non zero energy splitting.
We demonstrate the dependence of the energy splitting of the distance and radius of the vortices And show that a further exponential energy splitting in dependence on the distance between the vortex and the bulk edge exists. We visualise the adiabatic evolution of the quasi-particle exchange in real space. For the non-adiabatic regime, the dependence of the fidelity on the exchange speed is computed numerically for two and four vortices.  We discuss the adiabatic and diadiabatic limits and show a theoretical lower limit for the angular velocity for a non zero energy splitting and two vortices. 
\end{abstract}
\vspace{3cm}
\begin{abstract}
Die vorliegende Masterarbeit demonstriert, dass der nicht-abelsche Austausch von Majorana-Moden in den Wirbeln des chiralen 2D-topologischen p-Wellen-Supraleiters in der Anwesenheit von einer Energie-Aufspaltung möglich ist.
Wir zeigen die Abhängigkeit der Energie-Aufspaltung der Majorana-Moden in Abhängigkeit des Wirbelpaar-Abstandes und seines Radius auf und weisen eine weitere Energie-Aufspaltung in Abhängigkeit des Abstands von Wirbel und Supraleiter-Oberfläche nach. Der Austausch wurde für den adiabatischen Fall im Ortsraum  visualisiert.
Für die nicht-adiabatische Zeitentwicklung wurde die Fidelity in Abhängigkeit der Austauschgeschwindigkeit für zwei und vier Wirbel numerisch bestimmt.
Sowohl der adiabatische als auch der diabatische Fall wurden analytisch diskutiert und für den Fall einer Energieaufspaltung wurde eine untere Grenze für den Austausch von zwei Wirbeln beschrieben.
\end{abstract}
\pagebreak
\end{document}

